\documentclass[10pt]{beamer}
\usetheme{metropolis}
\usepackage{graphicx}
\usepackage{listings}
\usepackage{booktabs}
\usepackage{xcolor}
\usepackage{hyperref}

\definecolor{codebg}{HTML}{F5F5F5}
\definecolor{codegreen}{HTML}{2E7D32}
\definecolor{codegray}{HTML}{757575}
\definecolor{codeblue}{HTML}{1565C0}

\lstdefinestyle{rstyle}{
  language=R, backgroundcolor=\color{codebg},
  basicstyle=\ttfamily\scriptsize,
  keywordstyle=\color{codeblue}\bfseries,
  stringstyle=\color{codegreen},
  commentstyle=\color{codegray}\itshape,
  breaklines=true, frame=single,
  rulecolor=\color{codegray!30}, showstringspaces=false, tabsize=2,
  morekeywords={library,ggplot,aes,geom_col,geom_text,coord_flip,
    theme_minimal,theme,element_text,element_blank,labs,fct_reorder,
    ggsave,mutate,filter,count,arrange,slice_max}
}
\lstdefinestyle{pythonstyle}{
  language=Python, backgroundcolor=\color{codebg},
  basicstyle=\ttfamily\scriptsize,
  keywordstyle=\color{codeblue}\bfseries,
  stringstyle=\color{codegreen},
  commentstyle=\color{codegray}\itshape,
  breaklines=true, frame=single,
  rulecolor=\color{codegray!30}, showstringspaces=false, tabsize=4,
  morekeywords={import,from,as,pd,plt,sns,np,fig,ax}
}

\title{The Design Process: From Data to Display}
\subtitle{Workshop 1 --- Module 6}
\author{Prof.Asoc.\ Endri Raco}
\institute{Data Visualization for Data Scientists \\ UNYT Tirana}
\date{2025/26}

\begin{document}

\begin{frame}
  \titlepage
\end{frame}

% ================================================================
\begin{frame}{Module Roadmap}
  \begin{enumerate}
    \item The visualization pipeline: acquire → represent → refine
    \item Cairo's visualization wheel and design trade-offs
    \item Munzner's nested model: why → what → how
    \item Audience analysis and chart-type selection
    \item Iterative refinement: from sketch to publication
  \end{enumerate}
  \begin{exampleblock}{Learning Outcome}
    You will be able to articulate a systematic design process that
    begins with the audience's question and ends with a polished,
    principled chart.
  \end{exampleblock}
\end{frame}

% ================================================================
\section{The Visualization Pipeline}
% ================================================================

\begin{frame}{Five Steps: Acquire → Parse → Filter → Mine → Represent}
  \begin{center}
    \includegraphics[width=0.95\textwidth]{figures_m06/design_process}
  \end{center}
  \begin{alertblock}{Pro-Tip}
    Most analysts jump directly to Step 5 (Represent). Steps 1--4 consume
    80\% of the effort but determine 80\% of the chart's value. A
    beautiful chart built on dirty data is a beautiful lie.
  \end{alertblock}
\end{frame}

% ================================================================
\begin{frame}{From Raw Data to Final Chart: Worked Example}
  \begin{center}
    \includegraphics[width=0.95\textwidth]{figures_m06/worked_example}
  \end{center}
  \begin{exampleblock}{Data Insight}
    The same four numbers pass through four stages: raw table, default
    chart, correctly encoded bars, and a polished chart with direct
    labels and an accent colour on the highest value. Each step adds
    interpretive value without adding data.
  \end{exampleblock}
\end{frame}

% ================================================================
\section{Design Frameworks}
% ================================================================

\begin{frame}{Cairo's Visualization Wheel}
  \begin{center}
    \includegraphics[width=0.6\textwidth]{figures_m06/cairo_wheel}
  \end{center}
  \vspace{2pt}
  Six continuums of design tension. Every visualization sits somewhere
  on each axis; there is no single ``correct'' position --- only
  positions that are \emph{appropriate for the audience and purpose}.
\end{frame}

% ================================================================
\begin{frame}{Cairo's Six Design Tensions}
  \centering
  \begin{tabular}{lp{3.5cm}p{3.5cm}}
    \toprule
    \textbf{Axis} & \textbf{Left Pole} & \textbf{Right Pole} \\
    \midrule
    1 & Abstraction (statistical) & Figuration (pictorial) \\
    2 & Functionality (analytical) & Decoration (aesthetic) \\
    3 & Density (data-rich) & Lightness (minimal) \\
    4 & Multidimensional & Unidimensional \\
    5 & Originality (novel form) & Familiarity (standard form) \\
    6 & Novelty (surprising) & Redundancy (reinforcing) \\
    \bottomrule
  \end{tabular}

  \vspace{6pt}
  \begin{alertblock}{Pro-Tip}
    For analytical audiences (analysts, scientists), lean toward
    abstraction, functionality, and density. For general audiences,
    lean toward figuration, decoration, and lightness.
  \end{alertblock}
\end{frame}

% ================================================================
\begin{frame}{Munzner's Nested Model}
  \begin{center}
    \includegraphics[width=0.85\textwidth]{figures_m06/munzner_nested}
  \end{center}
\end{frame}

% ================================================================
\begin{frame}{Nested Model: Threats at Each Layer}
  \begin{enumerate}
    \item \textbf{Domain layer}: solving the \emph{wrong problem} ---
          the chart answers a question nobody asked
    \item \textbf{Abstraction layer}: wrong data type mapping ---
          treating ordinal data as nominal, or continuous as discrete
    \item \textbf{Encoding layer}: ineffective visual idiom ---
          using a pie chart for 20 categories, or area for precise
          comparisons
    \item \textbf{Algorithm layer}: slow rendering for large data ---
          plotting 10M points without aggregation or WebGL
  \end{enumerate}

  \vspace{4pt}
  \begin{exampleblock}{Data Insight}
    Errors at outer layers propagate inward. No amount of visual polish
    (layer 3) can fix a chart that answers the wrong question (layer 1).
    \textbf{Always validate outward-in.}
  \end{exampleblock}
\end{frame}

% ================================================================
\begin{frame}{Task Taxonomy: Why $\times$ What $\times$ How}
  \begin{center}
    \includegraphics[width=0.85\textwidth]{figures_m06/task_taxonomy}
  \end{center}
\end{frame}

% ================================================================
\section{Audience \& Chart Selection}
% ================================================================

\begin{frame}{Know Your Audience}
  \begin{center}
    \includegraphics[width=0.92\textwidth]{figures_m06/audience}
  \end{center}
\end{frame}

% ================================================================
\begin{frame}{Audience Drives Every Design Decision}
  \centering
  \begin{tabular}{lp{3cm}p{3cm}}
    \toprule
    \textbf{Decision} & \textbf{Executive} & \textbf{Analyst} \\
    \midrule
    Detail level  & 1 key metric   & Full dataset \\
    Interactivity & Static, annotated & Filters, drill-down \\
    Palette       & Grey + 1 accent  & Multi-hue qualitative \\
    Chart count   & 1 per slide      & Dashboard grid \\
    Annotation    & Bold callout     & Reference lines \\
    \bottomrule
  \end{tabular}

  \vspace{6pt}
  \begin{alertblock}{Pro-Tip}
    Ask three questions before opening your tool:
    (1) Who will read this?
    (2) What decision will they make?
    (3) How much time will they spend?
  \end{alertblock}
\end{frame}

% ================================================================
\begin{frame}{Chart Selection Matrix}
  \begin{center}
    \includegraphics[width=0.85\textwidth]{figures_m06/chart_decision}
  \end{center}
\end{frame}

% ================================================================
\section{Iterative Refinement}
% ================================================================

\begin{frame}{Sketch → Draft → Polished}
  \begin{center}
    \includegraphics[width=0.95\textwidth]{figures_m06/iteration}
  \end{center}
  \begin{exampleblock}{Data Insight}
    Iteration is not optional. Even Tufte's exemplars went through
    multiple drafts. Start rough (sketch the structure), then refine the
    encoding (draft), then polish the typography and colour (final).
  \end{exampleblock}
\end{frame}

% ================================================================
\begin{frame}[fragile]{Iterative Refinement in R}
  \begin{columns}[T]
    \begin{column}{0.52\textwidth}
\begin{lstlisting}[style=rstyle]
library(tidyverse)

# Step 1: Quick sketch
ggplot(df, aes(x = cat, y = val)) +
  geom_col()  # default everything

# Step 2: Correct form
ggplot(df,
  aes(x = fct_reorder(cat, val),
      y = val)) +
  geom_col(fill = "#64B5F6",
           width = 0.6) +
  coord_flip()

# Step 3: Polished
ggplot(df,
  aes(x = fct_reorder(cat, val),
      y = val)) +
  geom_col(fill = "#1565C0",
           width = 0.5) +
  geom_text(aes(label = val),
    hjust = -0.2, size = 3,
    fontface = "bold") +
  coord_flip() +
  theme_minimal() +
  theme(panel.grid.major.y =
    element_blank()) +
  labs(x = NULL, y = "Count",
    title = "Top Categories")
\end{lstlisting}
    \end{column}
    \begin{column}{0.45\textwidth}
      \includegraphics[width=\textwidth]{figures_m06/r_iteration}

      \vspace{4pt}
      \textbf{Three commits:}
      \begin{enumerate}
        \item \texttt{geom\_col()} --- see the data
        \item \texttt{coord\_flip()} + colour ---
              correct the form
        \item Theme + labels + sort ---
              publication quality
      \end{enumerate}
    \end{column}
  \end{columns}
\end{frame}

% ================================================================
\begin{frame}[fragile]{Iterative Refinement in Python}
  \begin{columns}[T]
    \begin{column}{0.52\textwidth}
\begin{lstlisting}[style=pythonstyle]
import matplotlib.pyplot as plt

cats = ["A","B","C","D","E"]
vals = [42, 58, 35, 67, 51]

# Step 1: Quick sketch
fig, ax = plt.subplots()
ax.bar(cats, vals)  # defaults

# Step 2: Correct form
fig, ax = plt.subplots()
idx = sorted(range(len(vals)),
    key=lambda i: vals[i])
ax.barh([cats[i] for i in idx],
    [vals[i] for i in idx],
    color="#64B5F6", height=0.6)

# Step 3: Polished
fig, ax = plt.subplots()
colors = ["#1565C0"] * 5
colors[-1] = "#E53935"  # max
ax.barh([cats[i] for i in idx],
    [vals[i] for i in idx],
    color=[colors[i] for i in idx],
    height=0.5)
for i, v in enumerate(
    [vals[j] for j in idx]):
    ax.text(v+0.8, i, str(v),
        va="center", fontsize=8,
        fontweight="bold")
\end{lstlisting}
    \end{column}
    \begin{column}{0.45\textwidth}
      \includegraphics[width=\textwidth]{figures_m06/py_iteration}

      \vspace{4pt}
      \textbf{Same three steps:}
      \begin{enumerate}
        \item \texttt{ax.bar()} --- default
        \item Sort + horizontal ---
              correct the form
        \item Accent colour + direct labels
              + spine removal
      \end{enumerate}
    \end{column}
  \end{columns}
\end{frame}

% ================================================================
\begin{frame}[fragile]{The Complete Pipeline in R (20 lines)}
  \begin{columns}[T]
    \begin{column}{0.55\textwidth}
\begin{lstlisting}[style=rstyle]
library(tidyverse)

# 1. ACQUIRE
df <- read_csv("googleplaystore.csv")

# 2-3. PARSE + FILTER
df_clean <- df |>
  filter(Type %in% c("Free","Paid")) |>
  count(Category, sort = TRUE) |>
  slice_max(n, n = 10) |>
  mutate(Category =
    fct_reorder(Category, n))

# 4-5. MINE + REPRESENT
ggplot(df_clean,
  aes(x = Category, y = n)) +
  geom_col(fill = "#1565C0",
           width = 0.6) +
  geom_text(aes(label = n),
    hjust = -0.2, size = 3) +
  coord_flip() +
  theme_minimal() +
  theme(panel.grid.major.y =
    element_blank()) +
  labs(title = "Top 10 App Categories",
    x = NULL, y = "Count",
    caption = "Source: Kaggle")
\end{lstlisting}
    \end{column}
    \begin{column}{0.42\textwidth}
      \textbf{Pipeline in 20 lines:}

      \vspace{4pt}
      Lines 1--3: Acquire (read CSV)

      Lines 5--10: Parse + Filter (clean
      types, count, top 10, reorder)

      Lines 12--23: Mine + Represent
      (bar chart with Tufte styling)

      \vspace{8pt}
      \begin{alertblock}{Pro-Tip}
        The pipe operator (\texttt{|>})
        makes the pipeline explicit:
        each verb is one step in the
        design process.
      \end{alertblock}
    \end{column}
  \end{columns}
\end{frame}

% ================================================================
\begin{frame}[fragile]{The Complete Pipeline in Python (25 lines)}
  \begin{columns}[T]
    \begin{column}{0.55\textwidth}
\begin{lstlisting}[style=pythonstyle]
import pandas as pd
import matplotlib.pyplot as plt

# 1. ACQUIRE
df = pd.read_csv(
    "googleplaystore.csv")

# 2-3. PARSE + FILTER
df_clean = (df
    .query("Type in ['Free','Paid']")
    .value_counts("Category")
    .head(10)
    .sort_values()
    .reset_index(name="n"))

# 4-5. MINE + REPRESENT
fig, ax = plt.subplots(figsize=(7,5))
ax.barh(df_clean["Category"],
    df_clean["n"],
    color="#1565C0", height=0.6)
for i, v in enumerate(df_clean["n"]):
    ax.text(v+5, i, str(v),
        va="center", fontsize=8)
ax.set_title("Top 10 App Categories",
    fontweight="bold")
ax.spines["top"].set_visible(False)
ax.spines["right"].set_visible(False)
plt.tight_layout()
plt.savefig("chart.pdf")
\end{lstlisting}
    \end{column}
    \begin{column}{0.42\textwidth}
      \textbf{Pipeline in 25 lines:}

      \vspace{4pt}
      Lines 1--5: Acquire

      Lines 7--12: Parse + Filter
      (method chaining mirrors the pipe)

      Lines 14--25: Mine + Represent

      \vspace{8pt}
      {\small Method chaining in pandas
      (\texttt{.query().value\_counts()
      .head().sort\_values()}) is the
      Python equivalent of R's pipe.}
    \end{column}
  \end{columns}
\end{frame}

% ================================================================
\section{Synthesis}
% ================================================================

\begin{frame}{The Design Process Checklist}
  Before creating any chart, complete this checklist:
  \begin{enumerate}
    \item \textbf{Who} is the audience? (Executive / Analyst / Public)
    \item \textbf{What} question does the chart answer?
    \item \textbf{What} data type? (Table / Network / Spatial / Temporal)
    \item \textbf{What} analytical goal? (Compare / Distribute / Relate /
          Compose / Trend)
    \item \textbf{How} --- choose the chart type from the decision matrix
    \item \textbf{Iterate} --- sketch $\rightarrow$ draft $\rightarrow$ polished
    \item \textbf{Validate} --- check against Munzner's four threats
  \end{enumerate}
\end{frame}

% ================================================================
\begin{frame}{Key Takeaways}
  \begin{enumerate}
    \item The \textbf{five-step pipeline} (acquire → parse → filter → mine
          → represent) ensures that data quality precedes visual design.
    \item \textbf{Cairo's wheel} frames design as a set of trade-offs, not
          a single optimum. The audience and purpose determine the position.
    \item \textbf{Munzner's nested model} identifies four layers of
          potential failure. Always validate outward-in.
    \item \textbf{Chart selection} follows from the analytical goal:
          comparison → bar; distribution → histogram; relationship → scatter.
    \item \textbf{Iteration} is mandatory: sketch (see the data), draft
          (correct the form), polished (publication quality).
  \end{enumerate}
\end{frame}

% ================================================================
\begin{frame}{Essential Readings for This Module}
  \begin{itemize}
    \item Cairo, A. (2012). \emph{The Functional Art}, Chapters 2--3
          (The Visualization Wheel, Choosing the Right Form).
    \item Munzner, T. (2014). \emph{Visualization Analysis and Design},
          Chapters 1--4 (Nested Model, Task Taxonomy).
    \item Kirk, A. (2019). \emph{Data Visualisation: A Handbook for
          Data Driven Design}, Chapters 4--6 (The Design Process).
    \item Schwabish, J. (2021). \emph{Better Data Visualizations},
          Chapter 2 (The Chart Chooser).
  \end{itemize}

  \vspace{6pt}
  \textbf{Next Module}: Critique \& Redesign Workshop (W01-M07)
\end{frame}

\end{document}
